\documentclass[11pt,letterpaper]{article}
\usepackage[utf8]{inputenc}
\usepackage[T1]{fontenc}
\usepackage[spanish,es-nodecimaldot,es-tabla]{babel}
\usepackage{mathptmx}
\usepackage{amsmath,amssymb}
\usepackage{graphicx}
\usepackage{booktabs,array}
\usepackage{tabularx}
\usepackage{float}
\usepackage{enumitem}
\usepackage[hyphens]{url}
\usepackage[hidelinks]{hyperref}
\usepackage{geometry}
\usepackage{xcolor}
\usepackage[most]{tcolorbox}
\usepackage{pdflscape}
\geometry{letterpaper,left=2.4cm,right=2.4cm,top=2.3cm,bottom=2.3cm}
\linespread{1.1}
\setlength{\parindent}{0pt}
\setlength{\parskip}{0.5em}
\setlist[itemize]{leftmargin=1.4em,topsep=1pt,itemsep=1pt}
\setlist[enumerate]{leftmargin=1.7em,topsep=1pt,itemsep=1pt}
\definecolor{iiBlue}{HTML}{1F4E79}
\definecolor{iiOrange}{HTML}{C55A11}
\newtcolorbox{propbox}[1]{colback=iiBlue!6,colframe=iiBlue,coltitle=white,fonttitle=\bfseries,title=#1,arc=2pt,boxrule=0.9pt,left=6pt,right=6pt,top=3pt,bottom=3pt}
\newtcolorbox{alertbox}[1]{colback=iiOrange!7,colframe=iiOrange,coltitle=white,fonttitle=\bfseries,title=#1,arc=2pt,boxrule=0.9pt,left=6pt,right=6pt,top=3pt,bottom=3pt}
\newcommand{\hrulethin}{\noindent\rule{\textwidth}{0.6pt}}
\begin{document}\pagestyle{plain}
\begin{center}
    \textbf{\large Universidad de Santiago de Chile}\\[0.1cm]
    Departamento de Ingeniería Industrial
\end{center}
\vspace{0.15cm}\hrulethin\vspace{0.35cm}
\begin{center}
    {\scshape Estadística Aplicada \quad$\cdot$\quad Caso 2 · Clasificación supervisada}\\[0.35cm]
    {\LARGE \textbf{Rúbrica de evaluación}}\\[0.15cm]
    {\large Proyecto Integrador: Clasificación Supervisada}\\[0.2cm]
    Predicción de una respuesta categórica
\end{center}
\vspace{0.2cm}\hrulethin\vspace{0.4cm}
El producto se evalúa con \textbf{dos notas independientes}: \textbf{Cátedra} (presentación, oratoria, rigor e informe de respaldo) y \textbf{Laboratorio} (el script). Cada instrumento suma 100 puntos y se convierte por separado según la regla de la asignatura. \textbf{Niveles por dimensión:} 1 ausente o erróneo; 2 parcial con omisiones relevantes; 3 mínimo funcional;
4 claro, coherente y reproducible; 5 completo, crítico y defendible individualmente.

\section*{Cátedra --- ponderación}
\begin{center}\small
\begin{tabular}{p{4.9cm} p{8.2cm} c}
\toprule
Dimensión & Qué se evalúa & Puntos \\
\midrule
Oratoria y dominio conceptual & Expone sin leer, explica en vez de recitar, maneja los conceptos y responde de forma individual. & 45 \\
Calidad de la presentación & Diseño claro, una idea por lámina, visuales legibles con título y comentario; narrativa coherente. & 35 \\
Rigor técnico y decisiones & Método correcto y justificado (validación, baseline, métricas, interpretación). & 15 \\
Informe de respaldo & Máximo 5 páginas, coherente con la presentación y el script; respalda las afirmaciones. & 5 \\
\midrule
\textbf{Total} & & \textbf{100} \\
\bottomrule
\end{tabular}\end{center}
\begin{landscape}
\thispagestyle{plain}
\section*{Descriptores --- Cátedra}
\footnotesize\renewcommand{\arraystretch}{1.25}
\begin{tabular}{p{2.4cm} p{3.15cm} p{3.15cm} p{3.15cm} p{3.15cm} p{3.15cm}}
\toprule
Dimensión & Nivel 1 (20\%) & Nivel 2 (40\%) & Nivel 3 (60\%) & Nivel 4 (80\%) & Nivel 5 (100\%) \\
\midrule
Oratoria y dominio conceptual (45) & Lee la mayor parte de la exposición o no logra explicar; no responde las preguntas o comete errores conceptuales graves. & Recita apoyándose en notas o láminas; responde de forma vaga o imprecisa; dominio conceptual parcial. & Expone con apoyo puntual; explica las ideas centrales y responde lo básico, con vacíos menores. & Expone con fluidez sin leer, explica en vez de recitar y responde con solvencia la mayoría de las preguntas. & Domina y conecta los conceptos con la evidencia y responde con precisión preguntas fuera de lo previsto, de forma individual. \\
Calidad de la presentación (35) & Láminas sobrecargadas o ilegibles, con volcados de código o texto; sin narrativa. & Diseño irregular, varias ideas por lámina o figuras sin título ni comentario; narrativa débil. & Diseño funcional, una idea por lámina en general y figuras legibles; narrativa comprensible. & Diseño claro y consistente, una idea por lámina, figuras con título, unidades y comentario; narrativa coherente. & Diseño sobresaliente que refuerza el mensaje; cada lámina aporta una conclusión visible y la historia se sigue sin esfuerzo. \\
Rigor técnico y decisiones (15) & Método incorrecto o sin baseline; conclusiones no sostenibles o causalidad injustificada. & Método incompleto; comparaciones sin incertidumbre o decisiones poco justificadas. & Método correcto en lo esencial (validación, baseline, métricas) con justificación limitada. & Decisiones justificadas y coherentes con lo presentado; métricas e interpretación adecuadas. & Además, discute supuestos, sensibilidad y límites, y descarta complejidad innecesaria con evidencia. \\
Informe de respaldo (5) & Ausente, excede el límite de páginas o contradice la presentación o el script. & Incompleto o meramente descriptivo; respalda débilmente las afirmaciones. & Cumple el formato ($\leq 5$ páginas) y respalda las afirmaciones principales. & Claro, trazable y coherente con la presentación y el script; tablas y figuras sustentan decisiones. & Síntesis precisa que permite reconstruir la evidencia y la recomendación. \\
\bottomrule
\end{tabular}

\vspace{0.25cm}
\footnotesize La nota de cada dimensión se asigna por el nivel alcanzado (1 a 5), escalado a su puntaje máximo.
\end{landscape}

\section*{Cátedra --- puntaje por nivel}
Cada nivel equivale a un \% de logro (Nivel $k = k\times 20\%$); el puntaje es el máximo por ese \%:
\begin{center}\small
\begin{tabular}{l ccccc c}
\toprule
Dimensión & N1 & N2 & N3 & N4 & N5 & Máx \\
\multicolumn{1}{r}{\footnotesize \% de logro} & 20 & 40 & 60 & 80 & 100 & \\
\midrule
Oratoria y dominio conceptual & 9 & 18 & 27 & 36 & 45 & 45 \\
Calidad de la presentación & 7 & 14 & 21 & 28 & 35 & 35 \\
Rigor técnico y decisiones & 3 & 6 & 9 & 12 & 15 & 15 \\
Informe de respaldo & 1 & 2 & 3 & 4 & 5 & 5 \\
\midrule
\textbf{Total} & 20 & 40 & 60 & 80 & 100 & \textbf{100} \\
\bottomrule
\end{tabular}\end{center}
\section*{Cátedra --- penalizaciones}
Descuentos \textbf{fijos} sobre el subtotal (100), acumulables.
\begin{center}\small
\begin{tabular}{p{11cm} c}
\toprule
Situación & Descuento \\
\midrule
Lee gran parte de la exposición (láminas o notas) & $-15$ \\
No maneja los conceptos al ser preguntado & $-15$ \\
Informe ausente o supera 5 páginas & $-5$ \\
Excede el tiempo (>10 min) o no participan todos los integrantes & $-5$ \\
\bottomrule
\end{tabular}\end{center}
\clearpage
\section*{Laboratorio --- ponderación}
Nota independiente del script.
\begin{center}\small
\begin{tabular}{p{4.9cm} p{8.2cm} c}
\toprule
Dimensión & Qué se evalúa & Puntos \\
\midrule
Reproducibilidad y ejecución & Ejecuta desde una sesión limpia y reproduce los resultados presentados. & 40 \\
Pipeline sin fuga (leakage) & Preparación, selección y remuestreo dentro de los folds; test intacto. & 35 \\
Legibilidad, semillas y coherencia & Código legible y con semillas; los resultados coinciden con la presentación. & 25 \\
\midrule
\textbf{Total} & & \textbf{100} \\
\bottomrule
\end{tabular}\end{center}
\begin{landscape}
\thispagestyle{plain}
\section*{Descriptores --- Laboratorio}
\footnotesize\renewcommand{\arraystretch}{1.25}
\begin{tabular}{p{2.4cm} p{3.15cm} p{3.15cm} p{3.15cm} p{3.15cm} p{3.15cm}}
\toprule
Dimensión & Nivel 1 (20\%) & Nivel 2 (40\%) & Nivel 3 (60\%) & Nivel 4 (80\%) & Nivel 5 (100\%) \\
\midrule
Reproducibilidad y ejecución (40) & No ejecuta o no reproduce lo presentado; requiere reescritura. & Ejecuta con intervención manual no documentada o reproduce sólo parcialmente. & Reproduce los resultados principales desde una sesión limpia siguiendo instrucciones básicas. & Reproduce todo con una ejecución única y documentada. & Además, automatiza la ejecución y verifica la reproducibilidad (semillas globales, comprobaciones). \\
Pipeline sin fuga (leakage) (35) & Fuga evidente: preprocesamiento global, o el test participa del ajuste o la selección. & Validación con errores u operaciones fuera de los folds que alteran la estimación. & Separación entrenamiento-validación-test correcta y validación cruzada cuando corresponde. & Toda transformación, selección y remuestreo ocurren dentro de cada fold; el test permanece intacto. & Además, incluye verificaciones que detectan fuga, inconsistencias de partición o fuga temporal. \\
Legibilidad, semillas y coherencia (25) & Código ilegible o sin semillas; no corresponde con lo presentado. & Nombres inconsistentes o documentación mínima; resultados que no coinciden con la presentación. & Código legible, con semillas y documentación básica; coincide en lo esencial con la presentación. & Modular y bien documentado; los resultados coinciden exactamente con la presentación y el informe. & Además, es parametrizable, con manejo de errores y registro de configuraciones. \\
\bottomrule
\end{tabular}

\vspace{0.25cm}
\footnotesize La nota de cada dimensión se asigna por el nivel alcanzado (1 a 5), escalado a su puntaje máximo.
\end{landscape}

\section*{Laboratorio --- puntaje por nivel}
\begin{center}\small
\begin{tabular}{l ccccc c}
\toprule
Dimensión & N1 & N2 & N3 & N4 & N5 & Máx \\
\multicolumn{1}{r}{\footnotesize \% de logro} & 20 & 40 & 60 & 80 & 100 & \\
\midrule
Reproducibilidad y ejecución & 8 & 16 & 24 & 32 & 40 & 40 \\
Pipeline sin fuga (leakage) & 7 & 14 & 21 & 28 & 35 & 35 \\
Legibilidad, semillas y coherencia & 5 & 10 & 15 & 20 & 25 & 25 \\
\midrule
\textbf{Total} & 20 & 40 & 60 & 80 & 100 & \textbf{100} \\
\bottomrule
\end{tabular}\end{center}
\section*{Laboratorio --- penalizaciones}
Descuentos \textbf{fijos} sobre el subtotal (100), acumulables.
\begin{center}\small
\begin{tabular}{p{11cm} c}
\toprule
Situación & Descuento \\
\midrule
Fuga de información (leakage) o uso del test para seleccionar & $-25$ \\
Requiere cambios manuales no documentados para ejecutar & $-10$ \\
\bottomrule
\end{tabular}\end{center}
\section*{Focos técnicos (conceptos a dominar)}
\begin{itemize}
\item Prevalencia: PR-AUC frente a ROC-AUC.
\item Desbalance sólo dentro de los folds.
\item Calibración: log-loss/Brier vs ranking.
\item Umbral o top-$k$ por costos y capacidad.
\item Matriz de confusión y volumen de acciones.
\item Fuga de información y test intacto.
\end{itemize}
\section*{Co-evaluación (pondera ambas notas)}
Al cierre, cada integrante evalúa a sus compañeros mediante una encuesta. Para cada estudiante se calcula el
promedio de los puntajes recibidos, expresado como porcentaje de logro. Se aplica por igual a \textbf{ambas notas}
(cátedra y laboratorio).
\begin{itemize}
\item Si ese promedio es $\geq 60\%$, las notas individuales son las del grupo (no se alteran).
\item Si es $< 60\%$, cada nota individual se pondera por ese promedio (nota individual $=$ nota del grupo $\times$ logro de co-evaluación).
\item Quien no responda la encuesta recibe la nota mínima en su co-evaluación, sin afectar a los demás.
\end{itemize}
Su objetivo es que todos trabajen.
\end{document}
